\section{Solved problems}\label{sec:solved}

For each of the five problems we state the problem as it appears in the notebook, the answer, and the idea of the
proof; the complete proofs are in Appendix~\ref{app:proofs}. For each we also record what the agent knew about the
literature status.

\subsection{Problem 21.32 (P.~J.~Cameron): is ``$G\cong H'$'' decidable?}
\begin{problem}[21.32]
Is the following problem decidable, and if so, what is its complexity? Given a finite group $G$, is there a finite
group $H$ such that the derived subgroup of $H$ is isomorphic to $G$?
\end{problem}
\begin{theorem}\label{thm:2132}
Let $G$ be a finite group, $e=2\exp Z(G)$ and $m=\lceil\log_2|G|\rceil$. If $G\cong H'$ for some finite group $H$,
then $G\cong (H^*)'$ for a finite group $H^*$ with $|H^*|\le |\Aut G|\cdot|Z(G)|^2\cdot e^{2m}$. Consequently the
problem is decidable, and it lies in nondeterministic quasi-polynomial time (a witness of order $|G|^{O(\log|G|)}$).
Moreover, if $G\cong H'$ then some subgroup $Q$ with $\Inn(G)\le Q\le\Aut(G)$ satisfies $Q'=\Inn(G)$; when
$Z(G)=1$ this condition is also sufficient.
\end{theorem}
\emph{Idea.} Given $H$ with $H'=G$, put $K=C_H(G)$. Then $H/K\le\Aut G$, $[H,K]\le G\cap K=Z(G)$, and $K$ has
class $\le 2$. Three reductions shrink $H$ without changing $H'$: (1) replace $H$ by the subgroup generated by
$G$ and $2m$ elements whose commutators generate $G$; (2) the map $k\mapsto k^e$ is an endomorphism of $K$ whose
image $C$ is central in $H$ and of bounded index; (3) the central extension $H$ of $\bar H=H/C$ by $C$ can be
replaced, using the universal coefficient theorem and a Baer-sum argument, by a central extension $H^*$ of $\bar H$
by $C\cap G$ with the same derived subgroup. The last step is the only delicate one: one shows that the derived
subgroup of a Baer sum $(E_1\times_{\bar H}A)/\Delta$ is isomorphic to $E_1'$ when the Hopf map of $A$ vanishes.
The necessary condition follows by passing to $H/Z(G)$. \emph{Literature.} The agent found no statement of this
result; the problem was posed in the 21st issue (2026). Examples: $S_3$ and $D_8$ are not derived subgroups
(known), $Q_8=\SL(2,3)'$, $A_4=S_4'$.

\subsection{Problem 16.60 (D.~MacHale): $T(G)\ge |S_\alpha|$}
\begin{problem}[16.60]
For a finite group $G$ let $T(G)=\sum_{\chi\in\Irr(G)}\chi(1)$ and, for $\alpha\in\Aut(G)$, $S_\alpha=\{g\in G:
\alpha(g)=g^{-1}\}$. Is $T(G)\ge |S_\alpha|$ for all $\alpha$?
\end{problem}
\begin{theorem}\label{thm:1660}
Yes. For every finite group $G$ and $\alpha\in\Aut G$, with $C=C_G(\alpha^2)$ and $\sigma=\alpha|_C$ (an
involutory automorphism of $C$),
\[ |S_\alpha| \;=\; \sum_{\theta\in\Irr(C)}\varepsilon_\sigma(\theta)\,\theta(1)\;\le\; T(C)\;\le\; T(G), \]
where $\varepsilon_\sigma(\theta)\in\{1,0,-1\}$ is the twisted Frobenius--Schur indicator.
\end{theorem}
\emph{Idea.} $S_\alpha\subseteq C_G(\alpha^2)$ and $\alpha$ restricts to an involutory automorphism of that
subgroup; the twisted Frobenius--Schur theorem (Kawanaka--Matsuyama~\cite{KM}) counts $S_\sigma$ by indicators, and
$T$ is monotone on subgroups by Frobenius reciprocity. \emph{Literature.} The two ingredients are classical; the
reduction to the centraliser of $\alpha^2$ is the only new step, and the agent regards this answer as ``possibly
known or easy for experts''. The classical case $\alpha=\mathrm{id}$ is the inequality $T(G)\ge 1+\#\{\text{involutions}\}$.

\subsection{Problems 2.78 / 3.57 (P.~I.~Trofimov): group-theoretic numbers}
\begin{problem}[2.78]
For a finite group $G$ let $f(G)$ be the number of orders $d$ such that $G$ has a non-normal subgroup of order $d$
(the number of $\mathrm{IE}\bar n$-systems). An integer $k$ is soluble if every group with $f(G)=k$ is soluble,
simple if $f(S)=k$ for some non-abelian simple $S$, absolutely simple if in addition no non-soluble non-simple
group has $f=k$, and composite if no simple group has $f=k$. Are the sets of soluble numbers and of absolutely
simple numbers finite or infinite? Do composite non-soluble numbers exist?
\end{problem}
\begin{theorem}\label{thm:278}
Both sets are finite: every $k\ge 86$ equals $f(A_5\times P)$ for a suitable $p$-group $P$ ($p\ge 7$), which is
non-soluble and not simple. Moreover $0,\dots,6$ and $8$ are soluble numbers, $7$ is absolutely simple
($f(G)=7$ iff $G\cong A_5$), and $f(S)\ge 9$ for every non-abelian simple $S\ne A_5$.
\end{theorem}
\emph{Idea.} A product formula $f(G\times H)=|\Ord(G)||\Ord(H)|-n(G)n(H)$ for coprime orders, and an explicit family
of $p$-groups $M_{p^n}\times C_{p^r}$ (modular group times cyclic) with $|\Ord|=n+r+1$ and $f=r+1$, give
$f(A_5\times P)=7n+9r+9$, which represents every integer $\ge 86$. The lower bound $f\ge 9$ for simple groups
other than $A_5$ uses Thompson's list of minimal simple groups and a count of subgroup orders; the case $f\in\{7,8\}$
is settled by an analysis of minimal normal subgroups. \emph{Literature.} The problem dates from the second issue
(1965) and is still listed as open; the existence of composite non-soluble numbers remains open (data suggests
that every $k\ge 9$ is a simple number).

\subsection{Problem 18.46 (Ant.~A.~Klyachko): sharpness of $|H|\le 2|G|^2$}
\begin{problem}[18.46]
Every finite group $G$ embeds in a finite group $H$ in which every element of $G$ is a square, and $H$ can be
chosen with $|H|\le 2|G|^2$. Is this estimate sharp? (It cannot be improved beyond $|G|^2$.)
\end{problem}
\begin{theorem}\label{thm:1846}
Yes. If $H\supseteq D\cong D_8$ and every element of $D$ is a square in $H$, then $|H|\ge 128=2|D_8|^2$; the wreath
product $D_8\wr C_2$ attains the bound.
\end{theorem}
\emph{Idea.} Exhaustive verification: all groups of order $8,16,\dots,120$ are in the SmallGroups library, and two
independent GAP computations show that none contains a dihedral subgroup of order 8 consisting of squares.
Exactly 14 of the 2328 groups of order 128 do. The table of minimal $|H|$ for all groups of order $\le 12$
shows that $D_8$ is the smallest group needing the factor 2. \emph{Literature.} The bound is from
Baranov--Klyachko (2012); the agent found no prior determination of the sharpness.

\subsection{Problem 13.19 (Yu.~M.~Gorchakov): regularity of $Q/H$}
\begin{problem}[13.19]
Suppose that a group $Q$ and its normal subgroup $H$ are subgroups of $G_1\times\dots\times G_n$ such that for each
$i$ the projections of both $Q$ and $H$ onto $G_i$ coincide with $G_i$. If $Q/H$ is a $p$-group, is $Q/H$ a
regular $p$-group?
\end{problem}
\begin{theorem}\label{thm:1319}
No, for every prime $p$. Every finite $p$-group $P$ of class $c$ occurs as $Q/H$ in such a configuration with
$n=c+1$. The smallest explicit example has $|Q|=2^9$, $n=3$ and $Q/H\cong D_8$.
\end{theorem}
\emph{Idea.} The reformulation is: $Q$ has normal subgroups $N_1,\dots,N_n$ with $\bigcap N_i=1$ and $HN_i=Q$
for all $i$ (then $G_i=Q/N_i$). The class of $Q/H$ is at most $n-1$, so non-regular quotients need $n\ge 3$. The
explicit example is a class-2 group of order $2^9$ built from three copies of $\F_2^2$ and three central
involutions; its properties were verified in GAP from a pc presentation. The general construction ``tags'' each
generator of $P$ by the subsets of $\{1,\dots,n\}$ containing a given index and uses a Möbius-inversion lemma on
word maps in nilpotent groups to show that $\bigcap N_i=1$. \emph{Literature.} The problem is from the 13th issue
(1995) and is listed as open; the class bound $\mathrm{cl}(Q/H)\le n-1$ is elementary.
